\begin{hint}\label{hint:sumapol2}
	Demuestre que la suma de
	los ángulos internos del polígono es
	igual a la suma de los ángulos
	internos de todos los triángulos.
\end{hint}

\begin{hint}\label{hint:sumapol1}
	Demuestre que es
	posible particionar el
	polígono en $n - 2$
	triángulos.
\end{hint}

\begin{hint}\label{elsalvador2015equilatero1}
	Considera qué pasa con el triángulo central.
\end{hint}

\begin{hint}\label{elsalvador2015equilatero2}
	Sin pérdida de generalidad dos fichas del triángulo
	central son rojas y la otra es azul.
\end{hint}

\begin{hint}\label{elsalvador2015equilatero3}
	Descartar casos.
\end{hint}

\begin{hint}\label{aliciabetomod3}
	Puede Beto asegurar que el número sea
	múltiplo de algo?
\end{hint}

\begin{hint}\label{aliciabetomod32}
	Múltiplos de $3$.
\end{hint}

\begin{hint}\label{ciclicoalgebra1}
	Teorema de pitágonas.
\end{hint}

\begin{hint}\label{ciclicoalgebra2}
	Teorema del coseno.
\end{hint}

\begin{hint}\label{ciclicoalgebra3}
	Construye un cuadrilátero cíclico.
\end{hint}

\begin{hint}\label{ciclicoalgebra4}
	Computa la razón con áreas.
\end{hint}

\begin{hint}\label{ricardojaja1}
	Lema \ref{lemma:righttrianglemindpoint}.
\end{hint}

\begin{hint}\label{rotomoteciacuadrilatero}
	Lema \ref{lemma:rotomotecia}
\end{hint}

\begin{hint}\label{telescopica1}
	$\frac1n - \frac1{n + 1} = \frac1{n(n + 1)}$.
\end{hint}
